\paragraph{核版 RH 临界五次域、分歧三次域与均匀性缺口的闭合接口}
在碰撞位势的核版 RH/Ramanujan 窗口、分歧晶格与 Fibonacci 共振常数已经建立之后，还可继续抽取出下列五条彼此衔接的算术--解析后果。它们分别刻画临界参数的数域次数、与分歧三次域的线性无交、RH 边界的纯谱性质、RH 窗口内部仍不可消除的均匀性缺口，以及最近分歧格点对高阶累积量增长与振荡的解析支配。

\begin{theorem}[核版 RH 临界参数生成五次域]\label{thm:xi-time-part9pb-rh-critical-quintic-field}
\leanverified{paper\_xi\_time\_part9pb\_rh\_critical\_quintic\_field}
沿用命题 \ref{prop:real-input-40-collision-rhk-window} 的记号，记
\[
Q(u):=u^5+4u^4+3u^3-96u^2-42u-14.
\]
则 \(Q\) 在 \(\QQ[u]\) 上不可约。因而若 \(u_R>1\) 为命题
\ref{prop:real-input-40-collision-rhk-window} 给出的唯一正实临界参数，则
\[
\boxed{\ [\QQ(u_R):\QQ]=5.\ }
\]
\end{theorem}
\begin{proof}
先把 \(Q\) 模 \(5\) 约化，得到
\[
\overline Q(u)=u^5-u^4-2u^3-u^2-2u+1\in \FF_5[u].
\]
对 \(a\in\FF_5=\{0,1,2,3,4\}\) 逐一代入，可知 \(\overline Q(a)\neq 0\)，故
\(\overline Q\) 没有一次因子。再直接检验 \(\FF_5[u]\) 中全部不可约二次因子，均不整除 \(\overline Q\)。由于五次多项式若在域上可约，则必出现 \(1+4\) 或 \(2+3\) 型分解，故 \(\overline Q\) 在 \(\FF_5[u]\) 上不可约。

又因 \(Q\) 为首一整系数多项式，Gauss 引理推出 \(Q\) 在 \(\QQ[u]\) 上亦不可约。命题 \ref{prop:real-input-40-collision-rhk-window} 已给出 \(u_R\) 满足 \(Q(u_R)=0\)，故 \(u_R\) 的极小多项式只能是 \(Q\) 本身，于是
\[
[\QQ(u_R):\QQ]=\deg Q=5.
\]
证毕。
\end{proof}

\begin{theorem}[临界五次域与分歧三次域线性无交]\label{thm:xi-time-part9pb-rh-branch-cubic-linear-disjoint}
沿用推论 \ref{cor:real-input-40-collision-branch-radius} 与定理
\ref{thm:xi-time-part9pb-rh-critical-quintic-field} 的记号。记
\[
\Delta(u):=4u^3+25u^2+88u+8.
\]
任取 \(\Delta(u)=0\) 的一个根 \(u_{\mathrm{br}}\)，并设
\[
K_{\mathrm{br}}:=\QQ(u_{\mathrm{br}}),
\qquad
K_{\mathrm{RH}}:=\QQ(u_R).
\]
则
\[
\boxed{\ K_{\mathrm{br}}\cap K_{\mathrm{RH}}=\QQ,\qquad
[K_{\mathrm{br}}K_{\mathrm{RH}}:\QQ]=15.\ }
\]
特别地，\(K_{\mathrm{br}}\) 与 \(K_{\mathrm{RH}}\) 在 \(\QQ\) 上线性无交。
\end{theorem}
\begin{proof}
先证 \(\Delta\) 在 \(\QQ[u]\) 上不可约。由有理根判别，其有理根只可能来自
\[
\pm 1,\ \pm 2,\ \pm 4,\ \pm 8,\ \pm \frac12,\ \pm \frac14,\ \pm \frac18.
\]
逐一代入均不为零，故 \(\Delta\) 没有有理根。由于 \(\deg \Delta=3\)，这已推出 \(\Delta\) 在 \(\QQ[u]\) 上不可约，从而
\[
[K_{\mathrm{br}}:\QQ]=3.
\]

另一方面，定理 \ref{thm:xi-time-part9pb-rh-critical-quintic-field} 已给出
\[
[K_{\mathrm{RH}}:\QQ]=5.
\]
若记
\[
E:=K_{\mathrm{br}}\cap K_{\mathrm{RH}},
\]
则塔式定理给出
\[
[E:\QQ]\mid [K_{\mathrm{br}}:\QQ]=3,
\qquad
[E:\QQ]\mid [K_{\mathrm{RH}}:\QQ]=5.
\]
因 \(3\) 与 \(5\) 互素，只能有
\[
[E:\QQ]=1,
\]
即
\[
K_{\mathrm{br}}\cap K_{\mathrm{RH}}=\QQ.
\]
于是
\[
[K_{\mathrm{br}}K_{\mathrm{RH}}:\QQ]
=
\frac{[K_{\mathrm{br}}:\QQ]\,[K_{\mathrm{RH}}:\QQ]}
{[K_{\mathrm{br}}\cap K_{\mathrm{RH}}:\QQ]}
=
\frac{3\cdot 5}{1}
=
15.
\]
证毕。
\end{proof}

\begin{theorem}[核版 RH 的临界边界是解析谱相变而非分歧碰撞]\label{thm:xi-time-part9pb-rh-transition-analytic-not-branch}
沿用推论 \ref{cor:real-input-40-collision-branch-radius} 与命题
\ref{prop:real-input-40-collision-rhk-window} 的记号。记
\[
\Delta(u)=4u^3+25u^2+88u+8,
\]
\[
\theta_R:=\log u_R,
\qquad
R_\theta:=\min_j\min_{k\in\ZZ}\bigl|\log u_j+2\pi i k\bigr|,
\]
其中 \(u_j\) 遍历 \(\Delta(u)=0\) 的三个根。则：
\begin{enumerate}
  \item 正实轴上不存在任何代数分歧点；
  \item 核版 RH 的临界点 \(u_R\) 位于 Perron 分支的解析域内部，而不在分歧边界上；
  \item 其到最近复分歧格点的解析裕量满足
  \[
  \boxed{\ R_\theta-\theta_R
  =
  R_\theta-\log u_R
  \approx 1.48342064735593>0.\ }
  \]
\end{enumerate}
因此，核版 RH 的临界边界来自解析域内部的谱支穿越，而不是由代数分歧碰撞触发。
\end{theorem}
\begin{proof}
对任意 \(u>0\)，多项式 \(\Delta(u)=4u^3+25u^2+88u+8\) 的各项都严格为正，故
\[
\Delta(u)>0\qquad (u>0).
\]
因此 \(\Delta(u)=0\) 的分歧点不可能落在正实轴上，第一条成立。

再由推论 \ref{cor:real-input-40-collision-branch-radius}，\(P_2(\theta)=\log\rho(e^\theta)\) 在 \(\theta=0\) 的 Taylor 收敛半径正是 \(R_\theta\)，而全部代数分歧点恰为
\[
\theta=\log u_j+2\pi i k,\qquad k\in\ZZ.
\]
另一方面，命题 \ref{prop:real-input-40-collision-rhk-window} 给出
\[
\textsf{RH}(u)\Longleftrightarrow 0<u\le u_R,
\qquad
u_R\approx 3.592621813443688,
\]
故
\[
\theta_R=\log u_R\approx 1.27888224610494.
\]
将其与
\[
R_\theta\approx 2.76230289346087
\]
比较，得到
\[
R_\theta-\theta_R\approx 1.48342064735593>0.
\]
于是 \(\theta_R\) 严格位于解析圆盘 \(\{|\theta|<R_\theta\}\) 的内部，从而第二条与第三条同时成立。

最后，在 \(u=u_R\) 处命题 \ref{prop:real-input-40-collision-rhk-window} 给出临界等号
\[
\Lambda(u_R)=\lambda_1(u_R)^{1/2}.
\]
故 RH 边界由正则谱支之间的不等式穿越所决定，而非由分歧方程 \(\Delta(u)=0\) 的塌缩所决定。证毕。
\end{proof}

\begin{theorem}[RH 窗口与常数级均匀性缺口严格共存]\label{thm:xi-time-part9pb-rh-window-uniform-gap-coexistence}
沿用命题 \ref{prop:real-input-40-collision-rhk-window}、推论
\ref{cor:fold-critical-resonance-gap} 与定理 \ref{thm:fold-critical-resonance-constant} 的记号。则在未加权点 \(u=1\) 有
\[
\textsf{RH}(1)\ \text{成立},
\]
但同时
\[
\boxed{\ 
\liminf_{m\to\infty}
F_{m+2}\left(\mathsf{Col}_m-\frac{1}{F_{m+2}}\right)
\ge
2C_\varphi^2
>0.\ }
\]
等价地，
\[
\boxed{\ 
\mathsf{Col}_m
\ge
\frac{1+2C_\varphi^2+o(1)}{F_{m+2}},
\qquad
1+2C_\varphi^2\approx 1.1007267225141177.\ }
\]
其中
\[
C_\varphi\approx 0.2244178274047294,
\qquad
2C_\varphi^2\approx 0.1007267225141179.
\]
因此，尽管 \(u=1\) 位于 kernel RH/Ramanujan 窗口内部，碰撞概率相对于均匀基线 \(1/F_{m+2}\) 仍保留一个不可消除的常数级缺口。
\end{theorem}
\begin{proof}
由命题 \ref{prop:real-input-40-collision-rhk-window}，
\[
\textsf{RH}(u)\Longleftrightarrow 0<u\le u_R,
\qquad
u_R>1.
\]
故 \(u=1\) 的确位于 RH 窗口内部，从而 \(\textsf{RH}(1)\) 成立。

另一方面，推论 \ref{cor:fold-critical-resonance-gap} 已给出
\[
\liminf_{m\to\infty}
F_{m+2}\left(\mathsf{Col}_m-\frac{1}{F_{m+2}}\right)
\ge
2C_\varphi^2,
\]
而定理 \ref{thm:fold-critical-resonance-constant} 给出 \(C_\varphi>0\)，故右端严格为正。将不等式改写，即得
\[
\mathsf{Col}_m
\ge
\frac{1+2C_\varphi^2+o(1)}{F_{m+2}}.
\]
代入现有数值近似
\[
C_\varphi\approx 0.2244178274047294
\]
便得到
\[
2C_\varphi^2\approx 0.1007267225141179,
\qquad
1+2C_\varphi^2\approx 1.1007267225141177.
\]
证毕。
\end{proof}

\begin{theorem}[最近分歧格点控制碰撞累积量的 Gevrey--1 增长与振荡渐近]\label{thm:xi-time-part9pb-branch-lattice-cumulant-gevrey-oscillation}
沿用推论 \ref{cor:real-input-40-collision-branch-radius} 与定理
\ref{thm:xi-time-part9pb-rh-transition-analytic-not-branch} 的记号。记 \(P_2(\theta)\) 为
\(\theta=0\) 处由 Perron 根定义的解析芽；在实轴上它与
\[
\log\rho(e^\theta)
\]
一致。再记
\[
\kappa_n:=P_2^{(n)}(0),
\qquad
R_\theta:=\min_j\min_{k\in\ZZ}\bigl|\log u_j+2\pi i k\bigr|.
\]
设最近分歧格点由一对共轭点
\[
\theta_\star,\ \overline{\theta_\star}
\]
实现，其中
\[
|\theta_\star|=R_\theta,
\qquad
\Im \theta_\star\neq 0,
\]
并且 Perron 解析支在这两个点处都呈简单平方根分歧。则：
\begin{enumerate}
  \item 存在常数 \(C>0\)，使得对一切 \(n\ge 1\) 都有
  \[
  \boxed{\;
  |\kappa_{2n}|
  \le
  C\,(2n)!\,R_\theta^{-2n}.\;}
  \]
  \item 存在复常数 \(c_\star\neq 0\)，使得当 \(n\to\infty\) 时，
  \[
  \boxed{\;
  \kappa_{2n}
  =
  2\Re\!\Bigl(
  c_\star\,(2n)!\,\theta_\star^{-2n}\,n^{-3/2}
  \Bigr)
  +O\!\bigl((2n)!\,R_\theta^{-2n}\,n^{-5/2}\bigr).\;}
  \]
\end{enumerate}
从而偶阶累积量的符号与幅度都呈现由 \(\arg\theta_\star\) 决定的确定性振荡；最近分歧晶格在实轴喷射上的干涉条纹由此直接外显。
\end{theorem}
\begin{proof}
由推论 \ref{cor:real-input-40-collision-branch-radius}，\(P_2\) 在圆盘
\[
|\theta|<R_\theta
\]
上全纯。于是对任意 \(0<r<R_\theta\)，由 Cauchy 估计可得
\[
|\kappa_{2n}|
=
\bigl|P_2^{(2n)}(0)\bigr|
\le
\frac{(2n)!}{r^{2n}}\sup_{|\theta|=r}|P_2(\theta)|.
\]
令 \(r\uparrow R_\theta\)，并把圆周上的上确界吸收到常数中，即得第一条 Gevrey--1 上界。

再看第二条。由简单平方根分歧假设，存在 \(P_2\) 在两处分歧点附近的局部展开
\[
P_2(\theta)
=
H_+(\theta)+b_+\Bigl(1-\frac{\theta}{\theta_\star}\Bigr)^{1/2},
\]
\[
P_2(\theta)
=
H_-(\theta)+b_-\Bigl(1-\frac{\theta}{\overline{\theta_\star}}\Bigr)^{1/2},
\]
其中 \(H_\pm\) 全纯，且 \(b_\pm\neq 0\)。由于这对共轭点是全部最近奇点，标准分支点 Darboux 传递给出 Taylor 系数渐近
\[
[\theta^N]P_2(\theta)
=
\frac{b_+}{\Gamma(-1/2)}\,N^{-3/2}\theta_\star^{-N}
+\frac{b_-}{\Gamma(-1/2)}\,N^{-3/2}\overline{\theta_\star}^{-N}
+O(R_\theta^{-N}N^{-5/2}).
\]
又因 \(P_2\) 在实轴上取实值，其 Taylor 系数皆为实数，故 \(b_-=\overline{b_+}\)。将 \(N=2n\) 代入并把因子 \(2^{-3/2}\) 吸收到常数中，令
\[
c_\star:=\frac{2^{-3/2}b_+}{\Gamma(-1/2)},
\]
再用
\[
\kappa_{2n}=(2n)!\,[\theta^{2n}]P_2(\theta),
\]
便得到
\[
\kappa_{2n}
=
2\Re\!\Bigl(
c_\star\,(2n)!\,\theta_\star^{-2n}\,n^{-3/2}
\Bigr)
+O\!\bigl((2n)!\,R_\theta^{-2n}\,n^{-5/2}\bigr).
\]
这就证明了第二条；振荡性则由 \(\theta_\star^{-2n}=R_\theta^{-2n}e^{-2in\arg\theta_\star}\) 直接给出。证毕。
\end{proof}

\noindent
由此，纯碰撞位势的 RH 临界参数不但生成一个独立的五次域，而且与分歧三次域在线性代数层面完全分离；同时，这一 RH 边界在解析几何上属于谱支穿越而非分歧塌缩，并且即便在最自然的未加权点 \(u=1\) 处，RH 窗口也并不迫使碰撞统计回到均匀基线。进一步，最近复分歧格点还把高阶累积量的 Gevrey--1 增长与振荡喷射直接锁定在同一解析半径上。于是算术次数、复分歧、局部喷射与宏观非均匀性四者在同一接口上获得了严格闭合。

\endinput
